\documentclass[10pt,a4paper]{amsart}
\usepackage[margin=19mm]{geometry}
\usepackage{amsmath,amssymb,mathtools}
\usepackage[scr=rsfs,cal=euler]{mathalfa}
\usepackage{xfrac}
\usepackage[unicode,hidelinks,bookmarks=false]{hyperref}
\newcommand{\ie}{i.e.\ }
\newcommand{\cf}{cf.\ }
\newcommand{\resp}{respectively\ }
\newcommand{\Subsection}{\subsection}
\providecommand{\nobreakdash}{\nobreak\hbox{-}}
\setlength{\emergencystretch}{2em}
\multlinegap0pt
\usepackage{bm}
\usepackage{bbm}
\usepackage{dsfont}
\usepackage{xspace}
\usepackage{cancel}
\usepackage{mathtools}
\usepackage{amsthm}
\makeatletter
\@ifundefined{theorem}{\newtheorem{theorem}{Theorem}}{}
\@ifundefined{lemma}{\newtheorem{lemma}[theorem]{Lemma}}{}
\@ifundefined{observation}{\newtheorem{observation}[theorem]{Observation}}{}
\makeatother
\newcommand{\B}{\mathcal{B}(\R)}
\newcommand{\KK}{\mathcal{K}}
\newcommand{\PP}{\mathbb{P}}
\newcommand{\R}{\mathbb{R}}
\newcommand{\eps}{\varepsilon}
\title[Conditional uncorrelation equals independence]{Conditional uncorrelation equals independence}
\author{Dawid Tar{\l}owski}
\date{Source: arXiv:2406.01849v4 (CC BY 4.0)}
\begin{document}
\maketitle

\noindent\emph{Source and licence.} Conditional uncorrelation equals independence. Version arXiv:2406.01849v4 (CC BY 4.0), distributed under the Creative Commons Attribution 4.0 International licence, as recorded in the arXiv licence metadata for this exact version and shown on the abstract page https://arxiv.org/abs/2406.01849v4. Full rights notice: licenses/arxiv-2406.01849-CC-BY-4.0.txt.\par\smallskip

\noindent\textit{Editorial scope.} Counted unit: the Theorem whose source label is \texttt{T1} in the article (source0.tex lines 362--370, source label \texttt{T1}), delivered together with the article's own complete original proof of it (source0.tex lines 373--434, one \texttt{proof} environment of 62 lines). No mathematics is written, repaired or re-derived: statement and proof are the article's own text line for line. Required context retained uncounted: RNCH (lines 240--242); L2 (lines 249--261); OL2 (lines 271--275); LRN (lines 320--324). Every definition, display and label of those lines is retained unchanged. Omitted from this package and not redistributed: 1--239, 243--248, 262--270, 276--319, 325--361, 371--372, 435--767. Nothing inside a retained excerpt is omitted or altered: every retained body is delivered exactly as the article's lines are written, so no omission receipt of any kind accompanies this package. Citations: the retained excerpts carry no citation command at all, so no bibliography entry of the article is transcribed and no reference is redistributed. Reference adaptation: none was needed and none was made; every reference inside the delivered unit resolves inside it, so no unresolved reference or citation is produced. Printed slips kept literal: no printed word, symbol, hypothesis or number of the counted statement or of its proof was altered. Layout and class: the article's own class is \texttt{sn-jnl} with packages including graphicx, multirow, amsmath, amssymb, amsfonts, amsthm, mathrsfs, appendix, xcolor, textcomp, manyfoot, booktabs, algorithm, algorithmicx; the wrapper is the lane's pinned \texttt{amsart} editorial wrapper at 10pt on A4, which restates the theorem-like environments the retained text needs and adds the article's own macro definitions that the retained text needs (\texttt{\textbackslash B}, \texttt{\textbackslash KK}, \texttt{\textbackslash PP}, \texttt{\textbackslash R}, \texttt{\textbackslash eps}), with extra wrapper packages (bm, bbm, dsfont, xspace, cancel, mathtools). The delivered numbering is the unit's own. Assets: no retained span contains a figure environment, a graphics-inclusion command, a PDF-page inclusion, a table or a TikZ picture; no image file is copied into this package, the package declares no asset dependency and no image file is redistributed with it. Licence: version arXiv:2406.01849v4 of this article is distributed under the Creative Commons Attribution 4.0 International licence, as recorded in the arXiv licence metadata for this exact version and shown on the abstract page https://arxiv.org/abs/2406.01849v4. A full rights notice is delivered at licenses/arxiv-2406.01849-CC-BY-4.0.txt. Characters: every retained excerpt is pure ASCII, so the delivered engine, XeLaTeX with its default Latin Modern text font, reports no missing character.
\begin{equation}\label{RNCH}
\frac{dQ}{dP}(x)=\lim\limits_{\eps\to0^+}\frac{Q(B(x,\eps))}{P(B(x,\eps))}\mbox{  \  \  \ \ \ for\ \  }P  \mbox{\ \ - almost any }\ x\in\R^n,
\end{equation}

\begin{lemma}\label{L2}
\begin{enumerate}
\item For any random variable $X\colon\Omega\to\R$, the function
\begin{equation}\label{C1} (a,b)\to E[X|X\in [a,b] ]
\end{equation}
uniquely determines the probability distribution $\PP_X$ (the domain of the function is the set of pairs $a<b$ with $\PP_X([a,b])>0$). 
\item For any random variable $X\colon\Omega\to\R$ with $E|X|<\infty$, the function
\begin{equation}\label{C2} 
t\to E[X|X\geq t]
\end{equation}
uniquely determines the probability distribution $\PP_X$ (the domain of the function is the set of $t\in\R$ with $\PP[X\geq t]>0$.)
\end{enumerate}
\end{lemma}

\begin{observation}\label{OL2}
   If $P_1$ and $P_2$ are two Borel probability measures on $\R$ (with\\ $\int\limits_{\R}|x|dP_i(x)<+\infty$ in case $\KK=\{[t,\infty)\colon t\in\R\}$) and $\KK^+_{P_2}\subset \KK^+_{P_1}$, then conditioning 
\begin{equation}\label{W}\frac{\int_A xdP_2(x)}{P_2(A)}=\frac{\int_A xdP_1(x)}{P_1(A)}, \ \ \ A\in \KK^+_{P_2},\end{equation}
forces $\KK^+_{P_1}=\KK^+_{P_2}$ and, by Lemma \ref{L2}, $P_2(A)=P_1(A)$ for any $A\in\B$.
\end{observation}

\begin{lemma}\label{LRN}
Let $X$ and $Z$ be random variables with $Z\geq0$, $E[Z]<\infty$ and $E[|XZ|]<\infty$. Define the nonnegative measure $Q_1$ and the signed measure $Q_2$, both on $(\R,\B)$, by: 
$$Q_1(A)=E[1_A(X)Z],\ A\in\B,\mbox{ and } Q_2(A)=E[1_A(X)XZ],\ A\in\B.$$
We have $$\frac{dQ_2}{dQ_1}(x)=x.$$
\end{lemma}

\begin{theorem}\label{T1}
Let  $\mathcal{K}=\{[a,b]\colon a<b\}$. For any random variables $X$ and $Y$, the following conditions are equivalent:
\begin{enumerate}
\item[1)] $X$ and $Y$ are independent,
\item[2)] for any $U=\{X\in A, Y\in B\}$ with $A, B\in \mathcal{K}$ and $\PP_{(X,Y)}(A\times B)>0$,$$E[XY|U]=E[X| U]\cdot E[Y|U].$$

\end{enumerate}
Under additional assumption: $E|X|<\infty$ and $E|Y|<\infty$, the  bounded intervals may be replaced with the half-bounded intervals $\{[t,\infty)\colon\ t\in \R\}$.
\end{theorem}

\begin{proof}
%We will prove the theorem under the assumption $\KK=\{[a,b]\}_{a<b}$ or $\KK=\{[t,\infty)\}_{t\in\R}$.\\
 The implication $"1) \Rightarrow 2)"$ is a rather straigthforward calculation. We will prove the second implication $"2)\Rightarrow 1)"$.\\
\textbf{Step 1.}  Assume  that $E|X|<\infty$ and $E|Y|<\infty$, and $\KK=\{[a,b]\}_{a<b}$ or $\KK=\{[t,\infty)\}_{t\in\R}$.\\
\textbf{1a.}  First we consider the case $\PP[Y>0]=1$. Fix $B\in\KK$ with  $\PP[Y\in B]>0$ which implies that $E[Y1_B(Y)]>0$. More, for any $A\in\KK$ we have  
$$E[1_A(X)1_B(Y)]>0\Leftrightarrow E[Y1_A(X)1_B(Y)]>0.$$

From condition 2) of the theorem, for any $A\in\KK$ with $\PP_{(X,Y)}(A\times B)>0$ we have

$$\frac{E[XY1_A(X)1_B(Y)]}{\PP_{(X,Y)}(A\times B)}=\frac{E[X1_A(X)1_B(Y)]}{\PP_{(X,Y)}(A\times B)}\cdot \frac{E[Y1_A(X)1_B(Y)]}{\PP_{(X,Y)}(A\times B)}$$
and thus
\begin{equation}\label{EEE}
\frac{E[XY1_A(X)1_B(Y)]}{E[Y1_A(X)1_B(Y)]}=\frac{E[X1_A(X)1_B(Y)]}{E[1_A(X)1_B(Y)]}.
\end{equation}
 Note that 
$$\hat{P_1}(A)=E[1_A(X)1_B(Y)]\mbox{ and }\hat{Q_1}(A)= E[Y1_A(X)1_B(Y)],\ A\in\B,$$ are nonnegative Borel measures with the same class of null sets $\hat{P_1}(A)=0\Leftrightarrow\hat{Q_1}(A)=0$. We now normalize those measures to obtain the corresponding  probability measures:
$$Q_1(A):=\frac{1}{E[Y1_B(Y)]}\cdot E[Y1_A(X)1_B(Y)],\ A\in \B$$
and
$$P_1(A):=\\P[X\in A| Y\in B]=\frac{E[1_A(X)1_B(Y)]}{E[1_B(Y)]},\ A\in \B.$$
 We want to reformulate  equation \eqref{EEE} in a way which will allow us to use Lemma \ref{LRN}. In this purpose define the following signed measures
$$Q_2(A):= \frac{1}{E[Y1_B(Y)]}\cdot{E[XY1_A(X)1_B(Y)]},\ \ A\in\B,$$
$$P_2(A):=\frac{1}{E[1_B(Y)]}\cdot E[X1_A(X)1_B(Y)], \ \ A\in\B,$$
so,  for any $A\in\KK$ with $\PP[X\in A,Y\in B]>0$, equation \eqref{EEE} takes the form
\begin{equation}\label{RN}
\frac{Q_2(A)}{Q_1(A)}=\frac{P_2(A)}{P_1(A)}.
\end{equation}
Now, we apply Lemma \ref{LRN}  to the pair $Q_2$ and $Q_1$ (with $Z:=Y1_B(Y)$)  which provides  the Radon Nikodym-derivative $\frac{dQ_2}{dQ_1}(x)=x$. In case $\KK=\{[a,b]\colon a<b\}$, Equation $\eqref{RN}$ forces $\frac{dP_2}{dP_1}=\frac{dQ_2}{dQ_1}$. Indeed, by  \eqref{RNCH},  for  $(P_1+Q_1)$ - almost any $x\in\R$,
$$\frac{dP_2}{dP_1}(x)=\lim\limits_{\eps\to0^+}\frac{P_2([x-\eps,x+\eps])}{P_1([x-\eps,x+\eps])}\mbox{ and }\frac{dQ_2}{dQ_1}(x)=\lim\limits_{\eps\to0^+}\frac{Q_2([x-\eps,x+\eps])}{Q_1([x-\eps,x+\eps])},$$
and the above limits are equal to each other by \eqref{RN}. However,  in case $\KK=\{[t,\infty)\colon t\in\R\}$ we need to use Lemma \ref{LRN} again:  the pair $P_1$ and $P_2$, by Lemma \ref{LRN} with $Z:=1_B(Y)$, satisfies  $\frac{dP_2}{dP_1}(x)=x$. Hence, in both cases, the equation \eqref{RN} takes the form 
$$\frac{\int\limits_AxQ_1(dx)}{Q_1(A)}=\frac{\int\limits_AxP_1(dx)}{P_1(A)},\ \mbox{ for any } A\in\KK\mbox{ with }P_1(A)>0.$$
By Lemma \ref{L2}, we have $Q_1(A)=P_1(A)$, $A\in \B$. As $B\in\KK$ with $\PP_Y(B)>0$ was fixed arbitrarily, by definitions of $Q_1$ and $P_1$, we have obtained that for any $B\in\KK$ and  $A\in\B$ with $\PP_{(X,Y)}(A\times B)>0$, we have 
$$\frac{E[Y1_A(X)1_B(Y)]}{E[1_A(X)1_B(Y)]}=\frac{E[Y1_B(Y)]}{E[1_B(Y)]}.$$

In other words, for $A\in\B, B\in\KK$  with $\PP[X\in A, Y\in B]>0$,
\begin{equation}\label{EXY}
E[Y|X\in A,\ Y\in B]=E[Y| Y \in B].
\end{equation}
 As the conditioning $E[Y|X\in A, Y\in B]$ does not depend on $X$, it is  natural  that $X$ and $Y$ are independent. To show this, we will use Lemma \ref{LRN} again: Equation \eqref{EXY} may be rewritten as:
$$\frac{\frac{1}{\PP[X\in A]}\cdot E[Y1_A(X)1_B(Y)]}{\PP[Y\in B| X \in A]}=\frac{E[Y1_B(Y)]}{\PP[Y\in B]},$$
 and hence, by Lemma \ref{LRN}, 
$$\frac{\int\limits_B y\PP_Y[ dy| X \in A]}{\PP[Y\in B| X \in A]}=\frac{\int\limits_By\PP_Y(dy)}{\PP_Y[ B]},$$
where $A\in\B,\ B\in\KK$  are arbitrary with $\PP_{(X,Y)}(A\times B)>0$.
%\mbox{ for any }B\in \KK\mbox{ with }\PP[Y\in B|X\in A]>0.$$
%which, for any $A\in \B$ with $\PP_X(A)>0$, 
 By Observation \ref{OL2}, we have $\PP_Y(\ \cdot\ |X\in A)=\PP_Y(\cdot)$, where  $A\in\B$ is arbitrary with $\PP[X\in A]>0$. This implies that $X$ and $Y$ are independent. 

\textbf{1b.} Now consider the case $Y>-c$, where $c>0$. By the assumption, $X$ and $Y$ are conditionally uncorrelated in sense $cov_{(A,B)}(X,Y)=0$ for any admissible $A$,$B\in \KK$. By \textbf{1a.}, it is enough to show that $X$ and $Y+c$ are conditionally uncorrelated too. Note that for any $A,B\in \KK$  we have $\PP_{(X,Y+c)}(
A\times B)>0\Leftrightarrow \PP_{(X,Y)}(A\times(B-c))>0$, where $B-c=\{b-c\colon b\in B\}$. Furthermore, the covariance is translation-invariant and the class $\KK$ is translation invariant in sense $B\in\KK \Rightarrow B-c\in\KK.$ Thus, for any $A,B\in\KK$ with $\PP_{(X,Y+c)}(A\times B)>0$, we have:
$$cov_{(A,B)}(X,Y+c)=cov(X,Y+c|X\in A, Y+c\in B)=$$
$$=cov(X,Y|X\in A, Y+c\in B)=cov(X,Y|X\in A, Y\in B-c)=0.$$
 Now, as we see that $Y+c$ and $X$ are conditionally uncorrelated, by point \textbf{1a.} $X$ and $Y+c$ are independent. Hence, $X$ and $Y$ are independent.

\textbf{1c.} Now we do not assume that $Y$ is bounded from below. The random variables $X$ and $Y$ are independent if and only if for any $c\in\R$ with $\PP[Y>c]>0$ 
the $X$ and $Y$ are independent with respect to the conditional probability $\PP_c=\PP[\cdot| Y>c]$, i.e. 
$$\PP[X\in A, Y\in B|Y>c]=\PP[X\in A| Y>c]\cdot \PP[Y\in B| Y>c], \ A,B\in\B.$$
For any $c$ small enough let $cov_{\PP_c}$ denote the covariance calculated on the probability space $(\Omega,\Sigma,\PP_c)$. Note that for $A,B\in \KK$, if $B\cap [c,\infty)\neq \varnothing$ then $B\cap [c,\infty)\in\KK$. Furthermore,
$$cov_{\PP_c}(X,Y|X\in A, Y\in B)= cov(X,Y|X\in A,Y\in B \cap [c,\infty)).$$
Hence, $X$ and $Y$ are conditionally uncorrelated with respect to the conditional probability measure $\PP_c=\PP[\cdot|Y\geq c]$ which, by \textbf{1b.}, implies that $X$ and $Y$ are independent with respect to $\PP_c$. As $c$ may be arbitrarily small, $X$ and $Y$ are independent. 

\textbf{Step 2.} It remains to consider the case $\KK=\{[a,b]\colon a<b\}$ without assuming the existence of the first moments of $X$ and $Y$. To prove that $X$ and $Y$ are independent it is sufficient to prove that they are independent with respect to the conditional probability
$$\PP_{c}=\PP[\ \cdot\ |(X,Y)\in [-c,c]^2]$$
for any $c>0$ large enough to  have $\PP[(X,Y)\in [-c,c]^2]>0$.  On the probability space $(\Omega,\Sigma,\PP_c)$ the $X$ and $Y$ are integrable and conditionally uncorrelated by the assumptions (note that if $A\in \KK$ and $A\cap[-c,c]\neq\varnothing$ then $A\cap[-c,c]\in\KK$), which by \textbf{Step 1} implies that $X$ and $Y$ are $\PP_c$ - independent for any large $c$. \end{proof}

\end{document}
